% GENERATED FILE. Edit canonical lessons, then run npm run generate:books.
% canonical-source-hash: fnv1a64:6e990149ea836829
% canonical-lessons: ES-C42-la

\chapter{La --- The Pronoun Carries What Left}
\label{ch:pronoun-la}
\hlchaptermodality{78}

\begin{chapteropening}
\textbf{By the end of this chapter:} \emph{I can replace a feminine noun with la, and explain that the pronoun carries a gender the sentence no longer states.}

Complete ES-C42-la: turn Hago la comida into La hago and say what decides between lo and la.
\end{chapteropening}

\section[la]{la — \textquotedblleft{}it,\textquotedblright{} for a feminine noun}

\label{lesson:ES-C42-la}



\begin{quote}
Say \textquotedblleft{}I have it,\textquotedblright{} about the book. (\emph{Lo tengo}.) Now guess the feminine version before you read on.
\end{quote}

\begin{grammarlens}[title={Grammar Lens: the form you probably guessed}]
\begin{quote}
\textbf{la} — \emph{it}, standing in for something feminine.
\end{quote}

\begin{quote}
\emph{Hago \textbf{la comida}.} — I make the meal. \emph{\textbf{La} hago.} — I make it.
\end{quote}

If you guessed \emph{la}, you were not lucky. You were doing what the shape of Spanish invites you to do: \emph{el} and \emph{la} mark the two genders everywhere else, so the object pronouns mark them the same way.

The pronoun agrees with the \textbf{noun it replaces} — not with you, and not with the verb. \emph{El libro} is masculine, so \emph{lo}. \emph{La comida} and \emph{la casa} are feminine, so \emph{la}.

That last point is worth pausing on. \emph{La} here is not agreeing with anything present in the sentence. The noun is \textbf{gone}; the pronoun is carrying its gender on its behalf. This is the first time in Spanish that a word has had to remember something that is no longer being said.
\end{grammarlens}

\begin{cousinweb}
Last chapter had a small mystery: why \emph{lo} and not \emph{el}, when both come from the same Latin word? Here that mystery quietly disappears.

The feminine object form in Latin was \textbf{illam}. Worn down as an article it gave \emph{la}. Worn down as a pronoun it gave — \emph{la}. Both routes arrived at the same place, because the vowel that got lost on the masculine side (\emph{elo $\to$ el}) was never at risk on the feminine side.

So the masculine pair split into two shapes and the feminine pair did not. Not an exception to memorise: a single sound change that happened to bite one word and not the other.
\end{cousinweb}

\subsection*{Your turn}
Say these aloud:

\begin{itemize}
\raggedright
  \item \textquotedblleft{}Hago la comida\textquotedblright{} then \textquotedblleft{}La hago\textquotedblright{}
  \item \textquotedblleft{}lo tengo\textquotedblright{} then \textquotedblleft{}la hago\textquotedblright{} — one syllable carries the whole difference
\end{itemize}

\begin{tcolorbox}[breakable,colback=teal!4,colframe=teal!35!black,title={Before you move on}]
Which pronoun replaces a feminine noun? (\textbf{La}.) What decides it — you, the verb, or the noun? (The \textbf{noun}, even though it is no longer there.) Next: the same slot, for the two people actually talking.
\end{tcolorbox}
